\begin{theorem}\label{night:P3K.10}
\StatementRecord{night:P3K.10}{false as stated; correction proved}{S6\_06\_BLIND: GPT-6 Astra (OpenAI, ChatGPT Pro mode, effort Pro), recorded re-derivation in the checkpoint dated 18 November 2026}{\cref{night:P1.11}, \cref{night:P3F.10}, \cref{night:P3F.21}, \cref{night:P3F.22}, \cref{night:P3F.24}, \cref{night:P3F.7}, \cref{night:P3F.9}, \cref{night:P3K.6}, \cref{night:P3K.7}}\par\noindent\textbf{Physical-line correction, 17 November 2026.} The displayed physical class \(25\) is refuted. The same positive primitive OS45 family has physical integer \(3\), Eells--Kuiper invariant \(3/28\), and class \(3\). Its source, line, cap and topology proof is retained. K68--69 prove the missing dual-line and sphere-gluing arrow. The following earlier assertion and its argument are retained as the located historical claim; their corrected current consequences are proved in K68--69.\medskip\par 
Let \(S\) be the minimal regular elliptic surface with section \(O\)
whose affine equation is

\[
 y^2=x^3-3f(t)x+2g(t),\qquad
 f=(t^2+1)^2+t,\qquad
 g=(t^2+1)^3+\frac32t(t^2+1)+\frac38.
\]

Put \(G=(t^2+1,\sqrt3/2)\), \(R=8G\),
\(M=\mathcal O_S(R-O)\), and \(Q=G\).  Choose any one of the four
nodal fibres as the simple zero of the linearization.  At the
ordinary fibre \(t=0\), perform the multiplicity-one logarithmic
gluing whose integral circle clutch is \(\widehat\gamma\).
The resulting construction of Engel's Remark 8.5 gives a proper
holomorphic family of compact complex manifolds diffeomorphic to
the standard \(S^6\), over a punctured scalar disc
\(0<|\lambda|<c\).

For the positive \(\lambda\)-loop, with the complex orientation and
the Eells--Kuiper convention of P3F.22, the return satisfies

\[
 A=-3,\qquad \mu(\Sigma_\Phi)=-\frac3{28},\qquad
 [\Phi]=25\in\mathbf Z/28.
\]

In particular this is a generator.  The positive finite parameter
pullbacks of this same family realize every oriented class in
\(\Theta_7\).  For every nonzero such class the mapping torus is
not diffeomorphic to \(S^1\times S^6\).
\end{theorem}

\begin{proof}
\begin{enumerate}[label=\arabic*.]
\item \par\medskip\noindent\textbf{The source and its section.}\enspace More generally put
\(R_0=t^2+a\),
\(f=R_0^2+bt\), and
\(g=R_0^3+\frac32btR_0+\frac38b^2\).  Exact expansion gives

\[
 g^2-f^3=
 \frac34b^2a(t^2+a)^2+
 \frac18b^3t^3+\frac98b^3at+\frac9{64}b^4.
\]

\item At \(a=b=1\), this is
\((48t^4+8t^3+96t^2+72t+57)/64\).  Its four roots are simple
and none is a root of \(f\).  At infinity the minimal coefficients
are units and the discriminant has order eight.  Thus the fibres
are exactly \(I_8+4I_1\).  The point \(G\) satisfies the equation
by direct substitution.

\item Rationality can be checked without identifying the surface by a
table entry.  Set \(u=x-(t^2+a)\).  The equation becomes

\[
 y^2=3(ut-b/2)^2+u^2(u+3a).
\]

\item On \(u\ne0\), the two rational functions

\[
 A_0=\frac{y-\sqrt3(ut-b/2)}u,
 \qquad B_0=\frac{y+\sqrt3(ut-b/2)}u
\]

satisfy \(A_0B_0=u+3a\).  Conversely,

\[
 u=A_0B_0-3a,\qquad
 t=\frac{B_0-A_0}{2\sqrt3}+\frac b{2u},\qquad
 y=\frac u2(A_0+B_0),\qquad x=t^2+a+u.
\]

\item These are inverse rational maps on nonempty open sets.  The smooth
minimal model is therefore a rational elliptic surface.

\item The section \(G\) is disjoint from \(O\).  At infinity, with
\(z=1/t\), its minimal coordinates are
\(X=1+az^2\), \(Y=(\sqrt3b/2)z^3\).
The two local nodal branches have orders three and five when
\(a=b=1\): expanding the square root of
\((1+az^2)^2+bz^3\) cancels the order-three term on exactly
one branch, whose next nonzero term has order five.  Their product
has order eight.  Hence \(G\) meets component three of \(I_8\),
up to reversing the cyclic numbering.  Its height is

\[
 h(G)=2-\frac{3(8-3)}8=\frac18.
\]

\item The root lattice has rank seven and determinant eight.  Shioda--Tate
therefore gives Mordell--Weil rank one.  The unimodular surface
lattice and the height discriminant formula give
\(h(G)=m^2|T|^2/8\), where \(m\) is its index in the free
Mordell--Weil group and \(T\) its torsion subgroup.  Thus
\(m=|T|=1\).  This identifies the actual generator and the
torsion-free \(I_8\) source, including the component marking.

\item The section \(R=8G\) meets the identity component at \(I_8\).
Its height is eight, so the same intersection formula gives
\(R\cdot O=3\).  Exact doubling locates all three intersections:
the reduced denominator of \(x(8G)\) is

\[
 (6t+1)^2(360t^2+282t+199)^2.
\]

\item The three roots are distinct and this denominator is coprime to
the discriminant.  The numerator is coprime to it as well, as
checked in the saved exact group-law calculation.  Consequently
all three intersections are ordinary and simple.  In particular
the Picard point of \(M\) is nonidentity at every \(I_1\) and
at the smooth identity component of \(I_8\).

\par\medskip\noindent\textbf{The actual marked local system.}\enspace We give a comparison
with the displayed source, rather than merely exhibit matrices
having the correct local spectra.  At \(a=0,b=1\) the source has
\(I_9+3I_1\).  Its \(j\)-map has degree twelve and ramification
profiles \(3^4\), \(2^6\), and \((9,1,1,1)\) over
\(0,1728,\infty\).  Indeed \(f=t(t^3+1)\) has four simple
roots, \(g=t^6+\frac32t^3+\frac38\) has six simple roots,
and \(g^2-f^3=(8t^3+9)/64\); these polynomials are pairwise
coprime.  Their indicated ramification already totals twenty-two,
so Riemann--Hurwitz leaves no further branch points.

\item The associated trivalent modular graph is unique.  It has four
vertices and six edges, with three monogon faces.  Each monogon
uses a loop at a vertex and consumes two of its three half-edges.
The three loops therefore occupy three distinct vertices, whose
remaining half-edges must join the fourth vertex.  This is the
three-star with a loop on each leaf.  Its cyclic ordering is unique
up to orientation-preserving graph isomorphism.  This proves
uniqueness of the permutation cover with these profiles by the
usual path-lifting classification of covers of the three-punctured
sphere.

\item For \(v=(p,q)^t\), write

\[
 D_v=\begin{pmatrix}1-pq&p^2\\-q^2&1+pq\end{pmatrix},
 \qquad T=D_{(1,0)}.
\]

\item The subgroup generated by
\(T^9,D_{(-3,1)},D_{(0,1)},D_{(3,1)}\) has the same cover:
its twelve-coset permutation table, exported by the program, has
exactly the three profiles above.  The products of the displayed
matrices equal the identity.  Thus the modular covering comparison
identifies this tuple with the actual \(I_9+3I_1\) source, up to
the explicitly permitted choice of marked vanishing paths and a
single integral symplectic basis.  The lift from projective
monodromy is fixed by the trace-two local \(I_n\) matrices.

\item Near infinity the deformation in \(a\) is analytically
\(xy=z^8(z+c a+O(a^2,az))\), with \(c\ne0\), after absorbing
units.  Hence it splits the \(I_9\) into an \(I_8\) and an
\(I_1\) with the same vanishing cycle.  Transport to \(a=1\)
along a path avoiding the finitely many other discriminant
collisions gives the actual five-point tuple

\[
 E_0=T^8,\quad E_1=T,\quad E_2=D_{(-3,1)},\quad
 E_3=D_{(0,1)},\quad E_4=D_{(3,1)}.
\]

\item Such transport is an isomorphism of the integral local systems
over the punctured surfaces and carries the section \(G\).
The path choice only changes the marking by the corresponding
Hurwitz moves.  No assertion about the topology depends on choosing
a particular branch of this transport: the program verifies all
four choices of the nodal zero below.

\item In the cohomology convention \((\psi,u,w,\delta)\) of P3F.21,
the weighted Jacobi lift for \(M=8G-O,Q=G\) is

\[
 J(E,k,z)=
 \begin{pmatrix}
 1&cj-ah&dj-bh&z\\
 0&a&b&8j\\
 0&c&d&8h\\
 0&0&0&1
 \end{pmatrix},
 \quad E=\begin{pmatrix}a&b\\c&d\end{pmatrix},
 \quad k=\binom jh.
\]

\item The component of \(G\) fixes \(k_0=(3,0)^t\).  An integral
change of Abel--Jacobi lift also sets \(k_4=0\).  Here is the
simultaneous gauge explicitly.  Adding the common integral lift
\((a,b)^t\) changes the first cocycle by
\((E_0-I)(a,b)^t=8b(1,0)^t\), so \(b\) fixes its component-three
representative.  The fourth coefficient changes by \(-a+3b\);
varying \(a\) then sets it to zero without changing the first.
The remaining
cocycles have the forms
\(k_1=r(1,0)^t\), \(k_2=j(-3,1)^t\),
\(k_3=h(0,1)^t\).  The nonzero-linearization local fillings
have rank-one nilpotent image; this forces their central entries
to be \(9,8r^2,8j^2,8h^2\), respectively.  At the scalar-zero
\(I_1\), the central entry is \(-1\).  The actual global product
must be the identity.  Its source columns give
\(j=2h\), \(r=-3h-3\), and its central entry gives

\[
 16(h+1)(4h+5)=0.
\]

\item The unique integral solution is \(h=-1\).  This proves the
actual section-lift identification; it is not a bounded search
for a convenient tuple.  The resulting matrices are

\[
\begin{aligned}
 T_0&=\begin{pmatrix}1&0&3&9\\0&1&8&24\\0&0&1&0\\0&0&0&1\end{pmatrix},
 &T_1&=\begin{pmatrix}1&0&0&0\\0&1&1&0\\0&0&1&0\\0&0&0&1\end{pmatrix},\\
 T_2&=\begin{pmatrix}1&2&6&32\\0&4&9&48\\0&-1&-2&-16\\0&0&0&1\end{pmatrix},
 &T_3&=\begin{pmatrix}1&1&0&8\\0&1&0&0\\0&-1&1&-8\\0&0&0&1\end{pmatrix},\\
 T_4&=\begin{pmatrix}1&0&0&-1\\0&-2&9&0\\0&-1&4&0\\0&0&0&1\end{pmatrix}.
\end{aligned}
\]

\item Moving the linearization zero to any of the other three \(I_1\)
fibres adds one to the central entry of \(T_4\) and subtracts
one at the chosen fibre.  The central matrices commute, so the
global product is still the identity.  Every integer calculation
below has been executed for all four choices.

\par\medskip\noindent\textbf{Line, fillings, and the ordinary logarithmic gluing.}\enspace
The line \(M\) has multidegree zero on every component of
\(I_8\).  Translation by \(Q\) therefore preserves its
multidegree.  The vertical line
\(\operatorname{Hom}(t_Q^*M,M)\) has zero intersection with
every fibre component and is a whole-fibre pullback.  Its degree,
computed by restriction to \(O\), is

\[
 (8G)\cdot O+1-(7G)\cdot O+(-G)\cdot O
 =3+1-3+0=1.
\]

\item Thus it is \(\mathcal O(F)\), and its section may have its
sole simple zero at any chosen nodal fibre.  Scaling the induced
lift by sufficiently small \(\lambda\) makes it uniformly
contracting.  On the complement of that fibre its quotient is a
smooth Hausdorff complex threefold, as follows directly from the
norm contraction and a compact annular fundamental domain.

\item At \(I_8\), the Tate and suspension periods, in a multidegree-zero
line frame, have valuation vectors \((8,0)\) and \((-3,0)\).
If they are denoted by \(h\) and \(g\), respectively, take

\[
 H=h^{-1}g^{-3},\qquad K=h^3g^8.
\]

\item The change has determinant one and their first valuations are
\(1,0\).  The second multiplier of \(K\) is \(\lambda^8\)
times a nonvanishing unit.  Thus the ray fan with consecutive
vertices \((j,1)\) is the actual primitive rank-one filling.
Equivalently it is the contracting quotient of the smooth source
line complement; \(Q\) cyclically permutes all eight components.
This equivalence retains the source elliptic map and its section
line.  The same primitive rank-one model applies at the other
three nonzero-linearization nodal fibres.  At the chosen zero
the two valuations are \((1,0),(0,1)\).  The square triangulation
has unimodular cones and gives the smooth rank-two cusp with
charts \(q=z_0z_1z_2\), its toric period identifications, and
proper reduced central fibre.  These are the specified local
models, with all valuation changes integral.

\item The base point \(p=0\) is ordinary.  On its smooth torus disc,
and on the universal logarithm cover of the scalar parameter,
let \(\gamma(z,v)\) be the holomorphic period of the primitive
suspension cycle. Glue a second copy of this smooth torus disc
by the following punctured-disc isomorphism, directed from the
cap copy to the ordinary copy:

\[
 \Psi_v(\zeta,z)=
 \left(\zeta+\frac{\log z}{2\pi i}\gamma(z,v),z\right),
 \qquad \lambda=e^{2\pi i v}.
\]

\item A logarithm branch change adds an actual lattice period, so this
is a single-valued holomorphic map on the torus family, with
inverse given by subtraction.  The cap is smooth, its central
fibre has multiplicity one, and its boundary clutch is exactly
\(\widehat\gamma\). Indeed a cap meridian with fixed cap coordinate
is carried to the ordinary base meridian plus \(\gamma\), the
\(m=1,n=1\) marking used below. This verifies the allowed logarithmic
transformation at an ordinary fibre even though its multiplicity
is one.  All pieces are proper over their closed base discs and
agree by the displayed holomorphic maps on collars; their finite
gluing is a compact Hausdorff complex threefold.

\par\medskip\noindent\textbf{Integral topology.}\enspace Put \(H_i=T_i^{-t}\) on the
homology lattice \(\Lambda\).  The images of \(H_i-I\) have
Smith factors \((1),(1),(1),(1),(1,1)\).  The primitive toric
collapses kill exactly these vanishing lattices.  We first fix
an actual global origin for the auxiliary untwisted family, before
the ordinary logarithmic transformation.  The restriction
\(M|_O\) has degree \(R\cdot O-O^2=4\).  Choose its section
with divisor \(4t_c\).  Its nonzero values over the other base
points give an origin in the quotient.  Near \(t_c\), its line
character is \(V=q^4\) times a holomorphic unit.  Subtracting
four gamma periods changes this to the nonvanishing bounded
character \(\lambda^{-4}e^{-8\pi iC}\) times a unit and changes
the source coordinate by the bounded translation \(-4\Gamma\).
It therefore extends in the dense torus of the cusp component.
This proves that the origin is global, including all five singular
fibres.  Every base meridian based at this origin contracts in
its local filling.

\item Van Kampen begins with the affine torus-bundle group over the
punctured sphere.  After these base meridians are killed, all
generators lie in the abelian fibre lattice, and the conjugation
relations impose the images of \(H_i-I\).  Their stacked relation
matrix has Smith factors \((1,1,1)\) and quotient \(\mathbf Z\)
measured by \(\psi\).  The ordinary log cap adds a vector with
\(\psi\)-value one.  The augmented Smith factors are
\((1,1,1,1)\), proving \(\pi_1=1\), with every affine origin
correction retained in the kernel already killed.  In the chosen
global-origin marking the complete relations are the abelian
fibre relations, \((H_i-I)\Lambda=0\) for \(0\leq i\leq4\),
and \(\widehat\gamma=0\).  Changing an origin only adds a vector
of \(\ker\psi\), which is precisely the primitive subgroup
generated by those five image lattices.

\item The complete cohomology calculation is integral.  For degree \(q\)
let \(W_i=\bigwedge^qT_i\), let \(L_i\) be the integer column
lattice of \(W_i-I\), and use its Hermite basis.  The marked
parabolic complex is

\[
 \mathbf Z^{\binom4q}\xrightarrow{d^0}
 \bigoplus_{i=0}^4 L_i\xrightarrow{d^1}
 \mathbf Z^{\binom4q},
\]

where \(d^0(a)=((W_i-I)a)_i\) and
\(d^1((b_i))=\sum_i(W_0\cdots W_{i-1})b_i\).
This is obtained by cutting the marked punctured sphere and
attaching each invariant stalk; the product relation gives
\(d^1d^0=0\).  The code exports every matrix in every degree.
For degrees one, two, and three, their ranks are respectively
\((3,3),(5,5),(3,3)\), and every nonzero Smith factor of
both differentials is one.  The middle groups vanish integrally;
the end groups are free of rank one.

\item There are no finite quotient caps in this example.  At each
primitive rank-one filling the collapse is the cellular pinching
of one circle in \(T^2\), crossed with the other elliptic factor;
it identifies cohomology with the full invariant lattice.  At
the rank-two filling the two valuation vectors and triangulation
identify the model integrally with the cusp of Engel's Lemma 7.4.
For an independent local chain verification, use the actual
dual-moment-cell quotient and blowup-collapse comparison of P1.11
in rank two, with \(G=3I_2-\mathbf1\mathbf1^t\).  Its translated
simplex counts are \((1,3,2)\), with twelve codimension-one
incidences.  For a simplex \(\sigma\), put
\(Q_\sigma=\mathbf Z^2/L_\sigma\).  The differential from
\(\bigoplus_{\dim\sigma=p}\bigwedge^q Q_\sigma\) is precisely
the signed integral quotient matrix (1.11.1), now with
\(0\leq p,q\leq2\).  The program
\path{code/part3_fourfold/rank_two_cusp.py} exports all these matrices
and all their Smith transformations.  Their total homology is
free of ranks \((1,2,4,2,1)\).  The collapse is the block sum
of \(\bigwedge^q(\mathbf Z^2\to Q_\sigma)\); its induced
homology matrices are onto with only unit Smith factors.  The
only nontrivial nonidentity block on the associated graded groups,
in total degree two, is

\[
 \begin{pmatrix}0&0&-1&1\\1&0&0&0\\0&-1&0&1\end{pmatrix}.
\]

\item Support leaves no possible higher differential.  The finite
filtration therefore makes the actual collapse onto and primitive
in every degree.  Its dual image is monodromy invariant and
primitive, with the same rank as the full invariant lattice, so
the two lattices agree.  This separately verifies every local
specialization assertion of the cited lemma needed here.
These exact local comparisons
give \(R^qf_*\mathbf Z=j_*\bigwedge^q\Lambda^*\) everywhere.

\item The surviving invariant generators, in degrees zero through four,
are

\[
 1,\quad\psi,\quad\omega=8uw+\psi\delta,
 \quad\eta_3=\psi uw,\quad\eta_4=\psi uw\delta.
\]

\item The corresponding coinvariant generators are
\(1,[\delta],\eta_2,[uw\delta],[\eta_4]\), with
\([\psi\delta]=8\eta_2\), \([uw]=\eta_2\), all other
degree-two basis classes zero.  The only nonzero Leray columns
are zero and two.  There is no unaccounted intrinsic cusp
differential.  Inversion in the untwisted torus family about the
global origin just constructed extends over every local filling:
on primitive nodal charts
\(X_i=Z/q^i\), \(Y_i=q^{i+1}/Z\), it is the actual holomorphic
map \((X_i,Y_i)\mapsto(Y_i,X_i)\) in chart \(-i-1\).
The other elliptic coordinate is inverted as well.  On the
rank-two square fan the lattice map \((v,h)\mapsto(-v,h)\)
sends each cone to its centrally inverted cone; its dual monomial
maps are holomorphic chart isomorphisms and agree on overlaps.
These are integral automorphisms of the specified fans.
The bounded origin characters only add a toric unit translation,
which extends too.  Thus this is an involution of the actual
untwisted compact family.  It acts by \((-1)^q\) on \(R^q\).
Naturality of Leray then makes every untwisted \(d_2\) satisfy
\(d_2=-d_2\).  Its targets are the free groups just computed,
so every one is zero integrally.

\item Replacing the ordinary cap changes only its boundary translation.
The cellular prism for this actual translation changes \(d_2\)
by contraction with its integral clutch.  Consequently the
translation prism of the ordinary
logarithmic clutch gives
\(d_2(a)=H[\iota_{\widehat\gamma}a]\), where \(H\) is the
positive base class.  The four contractions are

\[
 [\iota_\gamma\psi]=1,\quad
 [\iota_\gamma\omega]=[\delta],\quad
 [\iota_\gamma\eta_3]=\eta_2,\quad
 [\iota_\gamma\eta_4]=[uw\delta].
\]

\item For an arbitrary integral affine origin correction, the same
calculation is exactly its \(\psi\)-coordinate; source and line
circle coordinates contribute zero to all four projections.
Thus every displayed differential is multiplication by one for
the actual glued family.  There are no higher differentials,
because the base has dimension two.  Only
\(H^0=H^6=\mathbf Z\) remains, without torsion or extension
ambiguity.  Together with simple connectivity this proves that
the total manifold is a homotopy six-sphere, and \(\Theta_6=0\)
identifies its smooth oriented type with the standard sphere.

\par\medskip\noindent\textbf{The closed positive parameter loop.}\enspace Write \(t_c\) for
the chosen nodal zero and put
\(h(t)=(t-t_c)/t\).  Since
\(\gamma(z,v+1)=\gamma(z,v)+\delta\), the smooth logarithmic
gluing changes by \((\log t)/(2\pi i)\delta\) under the
positive parameter deck.  The global ordinary deck is
\(-\log h(t)\delta/(2\pi i)\): near \(t=0\) it has exactly
that value up to a holomorphic unit logarithm, which extends on
the cap. More explicitly, write
\[
 D_{\rm ord}=-\frac{\log h(t)}{2\pi i}\delta,\qquad
 D_{\rm cap}=-\frac{\log(t-t_c)}{2\pi i}\delta.
\]
The second logarithm is holomorphic on the chosen smooth cap.
The directed gluing satisfies the exact compatibility
\[
 D_{\rm ord}\circ\Psi_v=\Psi_{v+1}\circ D_{\rm cap},
 \qquad
 D_{\rm ord}=D_{\rm cap}+\frac{\log t}{2\pi i}\delta .
\]
Here the \(D\)'s denote the respective translations from \(v\)
to \(v+1\). At the cusp use \(q=h(t)\), \(s=\log q/(2\pi i)\).
Then the actual periods are

\[
 \gamma(s,v)=\binom{\Gamma(q)}{C(q)+s+v},\qquad
 \delta=\binom01,
\]

with \(\Gamma,C\) holomorphic at zero, and the deck translates
by \(-s\delta\).  Its line-character multiplier \(q^{-1}\)
extends through the cusp: compose an inverse gamma lattice shift
with the bounded source translation \(\Gamma\) and the
nonvanishing multiplier \(\lambda e^{2\pi iC}\).  At every
other original fibre \(h\) is a unit.  This proves that the
glued family descends holomorphically over the punctured
\(\lambda\)-disc, with the stated positive closing map.

\item Here \(\gamma\) is the inward period of the actual contracting
suspension. In the physical Abel and line coordinates it is
\[
 B(a,V)=
 \bigl(a+\Gamma(q),\,e^{2\pi iC(q)}q\lambda V\bigr),
 \qquad \Gamma(q)=\log_{\rm Tate}(-G).
\]
For small \(q,\lambda\), its positive real direction decreases
\(\log|V|\). This fixes its orientation relative to the source
page. With \(\delta\) the positive phase of \(V\), the complex
orientation gives
\[
 \operatorname{or}(Y)
   =-\operatorname{or}(S)\,d\log|V|
   =\operatorname{or}(S)\,d\gamma .
\]
Thus the positive gamma in the Hopf page below is this inward
period. The normalized smooth norm coordinate increases in
this direction.

\item For a completely marked angular comparison use the auxiliary
logarithmic cover \(q=Q_0^2\), and the two origins

\[
 o=\binom{-\Gamma s}{-Cs-s^2/2-sv},\qquad
 \sigma=\binom{\Gamma v}{Cv+(v^2-v)/2}.
\]

\item The first has parameter difference \(-s\delta\); the second
has difference \(-s\delta+\gamma\).  Under \(s\mapsto s+2\),
their differences are \(-2\gamma-2\delta\) and zero.  Thus
they descend on this auxiliary cover.  In contrast, the ordinary
origin has one-step difference \(-\gamma-\delta/2\); we do
not impose a false one-step periodicity.  The cusp origin extends
by its bounded unit characters.  Subtracting the ordinary origin
removes the parameter deck and gives the one-node parameter chain
\(X_i=V/\lambda^i\), \(Y_i=\lambda^{i+1}/V\),
\(X_iY_i=\lambda\), over every nonzero \(q\).  Its origin
extends.  These are the two extending collar origins.  No proper
two-variable corner or descent of such a corner is used.

\item Their exact difference is

\[
 F=\binom{\Gamma(s+v)}{C(s+v)+(s+v)^2/2-v/2},
\]

and hence

\[
 F(s,v+\theta)-F(s,v)
 =\theta\gamma(s,v)+\frac{\theta(\theta-1)}2\delta.
\]

\item The additional line gauge is periodic, equals one at both
endpoints, and vanishes in the circle quotient.  The double cusp
deck changes this angular difference by \(2\theta\delta\),
so the retained line gauge has winding two.  In particular the
actual positive \(\lambda\)-loop makes one positive gamma turn,
with no parameter cover.  Every identity is executed in
\path{code/part3_fourfold/os45_origins.py}.

\item More explicitly, the ordinary and cusp coordinates are
\(\zeta-o\) and \(\zeta-\sigma\), so the cusp-to-ordinary
transition is translation by \(F=\sigma-o\). Its angular
increment is the positive translation displayed above. This
identifies the moving fixed sphere in the constant orbit with
the positive inward gamma motion used by the real trace.

\par\medskip\noindent\textbf{The actual page and the return index.}\enspace The line circle is
free away from the selected rank-two cusp.  Indeed its equality
with an iterate of the contracting lift forces that iterate to
be zero by the norm inequality, and then forces the circle
element to be the identity.  This includes the whole \(I_8\)
neighbourhood, where translation permutes the components.  The
ordinary logarithmic cap is free as well.  At the single cusp
the primitive fan gives the one fixed rational curve with normal
weights \((-1,+1)\).  The orbit-pair argument P3F.7 therefore
applies to this proved sphere: its quotient is the standard
\((S^5,S^2)\), and its Euler line has positive meridian one.

\item The gamma circle over the base has the one ordinary clutch of
degree one, so its total space is the standard Hopf \(S^3\).
For precision, in the two standard Hopf charts the normalized
sections \((1,t)/\sqrt{1+|t|^2}\) and
\((1/t,1)/\sqrt{1+|1/t|^2}\) have angular coordinates related
by \(\theta_\infty=\theta_0+\arg t\), exactly the clutch
above.  A regular Hopf binding has a complex-oriented disc page
whose projection to \(\mathbf P^1\) has degree one.  Its inverse
image in the actual orbit is the elliptic page \(V\), with
boundary \(T^3\); its regular cap is the displayed rational
elliptic surface \(S\).  The affine origin obstruction on the
disc with marked singular values vanishes by the free-meridian
cochain argument of P3F.10.  Thus this comparison is with the
actual page, including the retained \(I_8\), rather than with
an unrelated elliptic surface.

\item The actual circle line restricts to \(M=8G-O\), up to a base
pullback, by its fibrewise Picard point.  At \(I_8\) both lines
have multidegree zero, so the same normalized fibrewise comparison
extends across its cycle of components.  On the base with boundary
the pullback is topologically trivial.  Therefore the actual
relative page invariants are

\[
 s_\delta=M^2=-2-2(3)=-8,
 \qquad p_{\rm rel}=3\sigma(S)=-24.
\]

\item Near the selected \(I_1\) the source translation is in the smooth
identity component.  The product coordinate of P3F.9 and the real
phase-preserving compression of P3K.6 consequently apply with
page degree \(d=1\).  The page phase can be chosen to avoid the
short gamma arc of that compression.  Its exact graph boundary,
integral Euler-line restriction, and oriented complement bordism
are therefore the ones proved in P3K.6–7.  The angular computation
above identifies this bordism with the actual positive scalar
return.  With the two orientation reversals of P3K.7 retained,

\[
 A=\frac{12s_\delta-p_{\rm rel}}{24}
   =\frac{-96+24}{24}=-3.
\]

\item The typed Hopf transfer and Eells--Kuiper normalization P3F.22 give
the asserted oriented class \(25\).  Pullback by
\(\lambda=z^n\) is a proper holomorphic family on the punctured
disc and its positive loop is exactly the \(n\)-fold iterate.
The classes \(-3n\pmod{28}\), for positive \(n\), exhaust
\(\mathbf Z/28\).  The mapping-torus obstruction P3F.24 applies
to each nonzero class and to both spin structures.  This proves
all the assertions of the statement.

\end{enumerate}
\end{proof}